\documentclass[10pt,a4paper]{amsart}
\usepackage[margin=19mm]{geometry}
\usepackage{amsmath,amssymb,mathtools}
\usepackage[scr=rsfs,cal=euler]{mathalfa}
\usepackage{xfrac}
\usepackage[unicode,hidelinks,bookmarks=false]{hyperref}
\newcommand{\ie}{i.e.\ }
\newcommand{\cf}{cf.\ }
\newcommand{\resp}{respectively\ }
\newcommand{\Subsection}{\subsection}
\providecommand{\nobreakdash}{\nobreak\hbox{-}}
\setlength{\emergencystretch}{2em}
\multlinegap0pt
\usepackage{mathtools}
\usepackage{amssymb}
\usepackage{float}
\usepackage{csquotes}
\usepackage[normalem]{ulem}
\newtheorem{lemma}{Lemma}
\newcommand{\SL}{\mathrm{SL}}
\newcommand{\tr}{\mathrm{tr}}
\title{Surjectivity of Engel Words on $\SL_2(\mathcal{O})$ and $\mathrm{PSL}_2(\mathcal{O}_2)$}
\author{Ayon Roy, Anupam Singh}
\date{Source: arXiv:2606.18880v1 (CC BY 4.0)}
\begin{document}
\maketitle

\noindent\emph{Source and licence.} Surjectivity of Engel Words on $\SL_2(\mathcal{O})$ and $\mathrm{PSL}_2(\mathcal{O}_2)$. Version arXiv:2606.18880v1 (CC BY 4.0), distributed by arXiv under the Creative Commons Attribution 4.0 International licence, http://creativecommons.org/licenses/by/4.0/, as recorded in the arXiv licence metadata for this exact version and shown on the abstract page https://arxiv.org/abs/2606.18880v1. Full licence notice: \texttt{licenses/arxiv-2606.18880-CC-BY-4.0.txt}.\par\smallskip

\noindent\textit{Editorial scope.} Counted unit: the article's lemma 2 (source0.tex lines 540--553). The article's complete original proof of it is delivered with it (source0.tex lines 554--672, one \texttt{proof} environment of 119 lines). No mathematics is written, repaired or re-derived: the statement and the proof are the article's own lines, in the article's own notation, and no retained line is substituted or re-flowed.  Retained uncounted as the source-local context the proof needs: lines 450--450.  Nothing was omitted from any retained span, so the package carries no omission record.  No retained line contains a citation, so the wrapper carries no bibliography and no bibliography entry.  No retained line contains a figure environment, an image-inclusion command or drawing code, so no image is delivered and the package declares no asset dependency.  Editorial formatting only, in addition: an AMS \texttt{amsart} preamble that restates the article's own declarations and macros used by the retained lines, namely the environments \texttt{lemma} on separate counters, together with the standard packages those lines need.  The retained environments print in this document as Lemma 1 in order of appearance, whereas the article prints the same items with the numbering of its own full document.  The complete native source of the article as distributed in the arXiv e-print for this exact version is bundled unchanged as \texttt{sources/203102/source0.tex}.
\section{The Trace Polynomial}{\label{tr}}

\begin{lemma}{\label{tr 2}}
Let $A_0 = \left (\begin{array}{cc} a  & b \\ c & d \end{array}\right)$ and $B_0=\left (\begin{array}{cc} e & n \\
m & w  \end{array}\right )$ are in $\SL_2(k)$ such that $\tr(\Bar{e}_r(A_0, B_0)) = \lambda_0$ for some $\lambda_0\in k\backslash\{2\}$. Let $r\geq 2$ and  $$f_r(x_1, x_2, x_3, x_4, y_1, y_2, y_3, y_4) = \tr(\Bar{e}_r(x,y))-\lambda_0$$ be a polynomial in $8$ variables over $k$. Consider matrices $x=\left (\begin{array}{cc} x_1  & x_2 \\
        x_3 & x_4
\end{array}\right)$ and $y=\left (\begin{array}{cc}
       y_1  & y_2 \\
        y_3 & y_4
\end{array}\right)$ in $\SL_2(k)$ such that $$\tr(\Bar{e}_r(x,y))-\lambda_0=\Tilde{\psi}_r(s_{r-1},v)$$ where $s_{r-1} = \tr(\Bar{e}_{r-1}(x,y))$ is also a function of $s,u$ and $v$; where $s=x_1+x_4$, $u=x_1y_1+x_2y_3+x_3y_2+x_4y_4$, $v=y_1+y_4$. 
    
Then, the matrix $\left[\begin{array}{cccccccc}
\frac{\partial f_r}{\partial x_1} & \frac{\partial f_r}{\partial x_2} & \frac{\partial f_r}{\partial x_3} & \frac{\partial f_r}{\partial x_4} & \frac{\partial f_r}{\partial y_1} & \frac{\partial f_r}{\partial y_2} & \frac{\partial f_r}{\partial y_3} & \frac{\partial f_r}{\partial y_4}
\end{array}\right]\Big|_{\alpha_0} = \mathbf{0}$ implies $\left[\begin{array}{ccc}
\frac{\partial\Tilde{\psi}_r}{\partial s_{r-1}} & \frac{\partial\Tilde{\psi}_r}{\partial v}  \end{array}\right]\Big |_{(\gamma_{r-1},v_0)}=\mathbf{0}$; where $\alpha_0=(a,b,c,d,e,n,m,w)$ and $\beta_0=(s_0,u_0,v_0)$, $\gamma_0=(s_0,u_0)$, $\beta_0=(\gamma_0,v_0)$ and $\gamma_j=s_j(\gamma_{j-1},v_0)$ for $1\leq j\leq r-1$. Moreover, $s_0 = tr(A_0) = a+d$, $u_0 = tr(A_0B_0) = ae+bm+cn+dw$ and $v_0=tr(B_0)=e+w$.
\end{lemma}

\begin{proof} 
\textbf{Step I:}
First, we prove that, for any fixed $r\geq 2$
\begin{eqnarray*}
\left[\begin{array}{ccc}
    \frac{\partial\Tilde{\psi}_r}{\partial s} & \frac{\partial\Tilde{\psi}_r}{\partial u} & \frac{\partial\Tilde{\psi}_r}{\partial v}  \end{array}\right]\Big |_{\beta_0}=\mathbf{0}\implies \left[\begin{array}{ccc}\frac{\partial\Tilde{\psi}_r}{\partial s_1} & \frac{\partial\Tilde{\psi}_r}{\partial v}  \end{array}\right]\Big |_{(\gamma_1,v_0)}=\mathbf{0} \\\implies \left[\begin{array}{ccc}
    \frac{\partial\Tilde{\psi}_r}{\partial s_2} & \frac{\partial\Tilde{\psi}_r}{\partial v}  \end{array}\right]\Big |_{(\gamma_2,v_0)}=\mathbf{0}\implies\cdots\implies \left[\begin{array}{ccc}
    \frac{\partial\Tilde{\psi}_r}{\partial s_{r-1}} & \frac{\partial\Tilde{\psi}_r}{\partial v}  \end{array}\right]\Big |_{(\gamma_{r-1},v_0)}=\mathbf{0} ; 
    \end{eqnarray*}
where $\gamma_0=(s_0,u_0)$, $\beta_0=(\gamma_0,v_0)$ and $\gamma_j=s_j(\gamma_{j-1},v_0)$ for $1\leq j\leq r-1$. Here, the chain rule, analogous to the partial derivatives, is applicable. As $s_1$ is a function of $s,u,v$ therefore the condition $\left[\begin{array}{cccccccc}
  \frac{\partial\Tilde{\psi}_r}{\partial s} & \frac{\partial\Tilde{\psi}_r}{\partial u} & \frac{\partial\Tilde{\psi}_r}{\partial v}\end{array}\right]\Big|_{\beta_0}=\mathbf{0}$ gives the following equations:
\begin{equation}{\label{lemsimp i}}
\frac{\partial\Tilde{\psi}_r}{\partial s}|_{\beta_0}=\frac{\partial\Tilde{\psi}_r}{\partial s_1}|_{(\gamma_1,v_0)}\frac{\partial s_1}{\partial s}|_{\beta_0}=0\implies \frac{\partial\Tilde{\psi}_r}{\partial s_1}|_{(\gamma_1,v_0)}(2s_0-u_0v_0)=0 \end{equation}
 \begin{equation}{\label{lemsimp ii}}
   \frac{\partial\Tilde{\psi}_r}{\partial u}|_{\beta_0}=\frac{\partial\Tilde{\psi}_r}{\partial s_1}|_{(\gamma_1,v_0)}\frac{\partial s_1}{\partial u}|_{\beta_0}=0\implies \frac{\partial\Tilde{\psi}_r}{\partial s_1}|_{(\gamma_1,v_0)}(2u_0-s_0v_0)=0 
  \end{equation}
  \begin{equation}{\label{lemsimp iii}}
   \frac{\partial\Tilde{\psi}_r}{\partial v}|_{\beta_0}=\frac{\partial\Tilde{\psi}_r}{\partial v}|_{(\gamma_1,v_0)}=0.
\end{equation}
Claim: The partial derivative $\frac{\partial\Tilde{\psi}_r}{\partial s_1}|_{(\gamma_1,v_0)}= 0$. We prove this claim by assuming contrarily as follows: If possible, let $\frac{\partial\Tilde{\psi}_r}{\partial s_1}|_{(\gamma_1,v_0)}\neq 0$. Then by Equation~(\ref{lemsimp i}) and Equation~(\ref{lemsimp ii}) we obtain $s_0(4-v_0^2)=0$. Therefore either $v_0=\pm2$ or $s_0=0$. 
Here we recall the trace formula introduced in Section~\ref{tr}, viz., $s_1=\tr(\bar e_1(x,y))=s^2+u^2+v^2-suv-2$ and $s_{m+1}=\tr(\bar e_{m+1}(x,y))=s_m^2+2v^2-s_mv^2-2$ for $m\geq 1$.
  
\noindent \textbf{Case I.} Let $s_0 = 0$. This immediately implies $u_0 = 0$, by Equation~(\ref{lemsimp ii}). Then the trace formula gives $\gamma_1=s_1(\gamma_0)=v_0^2-2$ which further implies $$s_2(\gamma_1,v_0)=(v_0^2-2)^2+2v_0^2-(v_0^2-2)v_0^2-2=2.$$   Continuing this process, inductively we obtain $\lambda_0=s_r(\gamma_{r-1},v_0)=2$; which is a contradiction.

\noindent \textbf{Case II.} When $v_0=\pm 2$ then using Equation~(\ref{lemsimp i}) and Equation~(\ref{lemsimp ii}), we obtain the Table~\ref{tab: values 5 points}.
 \begin{table}[ht]
\scriptsize
\centering
\renewcommand{\arraystretch}{1.8}\begin{tabular}{|l|l|l|l|l|}
\hline
$s_0$ & $u_0$ & $v_0$ & $s_1$ & $\lambda_0$ \\ \hline
$u_0$   & $u_0$  & 2 & 2 & 2 \\ \hline
$-u_0$   & $u_0$   & -2 & 2 & 2 \\ \hline
\end{tabular}
\caption{Values of $(s_0,u_0,v_0,s_1,\lambda_0)$}
\label{tab: values 5 points}
\end{table}
In this case, by a similar argument to case I, we obtain $\lambda_0=2$, which again contradicts the assumption. Therefore, $\frac{\partial\Tilde{\psi}_r}{\partial s_1}|_{(\gamma_1,v_0)}=0$. Hence, the claim is true.

This claim and the Equation~(\ref{lemsimp iii}) together implies $\left[\begin{array}{ccc}
    \frac{\partial\Tilde{\psi}_r}{\partial s_1} & \frac{\partial\Tilde{\psi}_r}{\partial v}  \end{array}\right]\Big |_{(\gamma_1,v_0)}=\mathbf{0}$.
   
    
\textbf{Step II:} Secondly we prove that for $2\leq j\leq r-1$; 
\begin{eqnarray*}
\left[\begin{array}{ccc}
\frac{\partial\Tilde{\psi}_r}{\partial s_{j-1}} & \frac{\partial\Tilde{\psi}_r}{\partial v}  \end{array}\right]\Big |_{(\gamma_{j-1},v_0)}=\mathbf{0}\  \text{implies}\  \left[\begin{array}{ccc}
    \frac{\partial\Tilde{\psi}_r}{\partial s_j} & \frac{\partial\Tilde{\psi}_r}{\partial v}  \end{array}\right]\Big |_{(\gamma_j,v_0)}=\mathbf{0}.
     \end{eqnarray*}
Now the condition $\left[\begin{array}{ccc}
    \frac{\partial\Tilde{\psi}_r}{\partial s_{j-1}} & \frac{\partial\Tilde{\psi}_r}{\partial v}  \end{array}\right]\Big |_{(\gamma_{j-1},v_0)}=\mathbf{0}$
    gives the following equations:
    \begin{equation}{\label{lemsimp iv}}
    \frac{\partial\Tilde{\psi}_r}{\partial s_{j-1}}\big |_{(\gamma_{j-1},v_0)}=\frac{\partial\Tilde{\psi}_r}{\partial s_{j}}\big |_{(\gamma_j,v_0)}\frac{\partial s_{j}}{\partial s_{j-1}}\big |_{(\gamma_{j-1},v_0)}=0\implies \frac{\partial\Tilde{\psi}_r}{\partial s_{j}}\big |_{(\gamma_j,v_0)}(2s_{j-1}(\gamma_{j-2},v_0)-v_0^2)=0
  \end{equation}
\begin{equation}{\label{lemsimp v}}
\frac{\partial\Tilde{\psi}_r}{\partial v}\big |_{(\gamma_{j-1},v_0)} =0\implies \frac{\partial\Tilde{\psi}_r}{\partial v}\big |_{(\gamma_{j},v_0)}=0\implies\frac{\partial\Tilde{\psi}_r}{\partial s_{j-1}}\big|_{(\gamma_{j-1},v_0)}\frac{\partial s_{j-1}}{\partial v}\big|_{(\gamma_{j-2},v_0)}+ 2v_0(2-s_{r-1}(\gamma_{r-2},v_0))=0.
\end{equation}
By Equation~(\ref{lemsimp iv}) and Equation~(\ref{lemsimp v}) we have $v_0(2-s_{r-1}(\gamma_{r-2},v_0))=0$. This implies $v=0$ or $s_{r-1}(\gamma_{r-2},v_0)=2$. If possible let  $s_{r-1}(\gamma_{r-2},v_0)=2$. Then $\lambda_0=s_r(\gamma_{r-1},v_0)=2$; which is a contradiction. Therefore $v_0=0$. 
 
If possible let $v_0^2=2s_{j-1}(\gamma_{j-2},v_0)$ in Equation~(\ref{lemsimp iv}). Then we immediately get $s_{j-1}(\gamma_{j-2},v_0)=0$. This implies $s_j(\gamma_{j-1},v_0)=-2$. As $j\leq r-1$, therefore $\lambda_0=s_r(\gamma_{r-1},v_0)=2$ which is again a contradiction. Therefore the only possibility is $v_0^2\neq 2s_{j-1}(\gamma_{j-2},v_0)$. This immediately implies $\frac{\partial\Tilde{\psi}_r}{\partial s_{j}}\big |_{(\gamma_j,v_0)}=0$. Now, considering this together with the first part, it is easy to see that 
\begin{eqnarray*}
\left[\begin{array}{ccc}
\frac{\partial\Tilde{\psi}_r}{\partial s} & \frac{\partial\Tilde{\psi}_r}{\partial u} & \frac{\partial\Tilde{\psi}_r}{\partial v}  \end{array}\right]\Big |_{\beta_0}=\mathbf{0}\implies \left[\begin{array}{ccc}\frac{\partial\Tilde{\psi}_r}{\partial s_1} & \frac{\partial\Tilde{\psi}_r}{\partial v}  \end{array}\right]\Big |_{(\gamma_1,v_0)}=\mathbf{0} \implies\cdots\implies \left[\begin{array}{ccc}
\frac{\partial\Tilde{\psi}_r}{\partial s_{r-1}} & \frac{\partial\Tilde{\psi}_r}{\partial v}  \end{array}\right]\Big |_{(\gamma_{r-1},v_0)}=\mathbf{0}.
\end{eqnarray*}

\textbf{Step III:} As the final step, we prove that 
\begin{eqnarray*}
\left[\begin{array}{cccccccc}
\frac{\partial f}{\partial x_1} & \frac{\partial f}{\partial x_2} & \frac{\partial f}{\partial x_3} & \frac{\partial f}{\partial x_4} & \frac{\partial f}{\partial y_1} & \frac{\partial f}{\partial y_2} & \frac{\partial f}{\partial y_3} & \frac{\partial f}{\partial y_4}
\end{array}\right]\Big|_{\alpha_0} = \mathbf{0}\  \text{implies}\  \left[\begin{array}{cccccccc}
\frac{\partial\Tilde{\psi}_r}{\partial s} & \frac{\partial\Tilde{\psi}_r}{\partial u} & \frac{\partial\Tilde{\psi}_r}{\partial v}\end{array}\right]\Big|_{\beta_0}=\mathbf{0}.
\end{eqnarray*}
As $\Tilde{\psi}_r(s,u,v)=f(x_1,x_2,x_3,x_4,y_1,y_2,y_3,y_4)$ and $s,u,v$ all are the polynomials in $x_1,x_2,x_3,x_4,y_1,y_2,y_3,y_4$ therefore 
\begin{eqnarray*}
\frac{\partial f}{\partial x_i}=\frac{\partial f}{\partial s}\frac{\partial s}{\partial x_i}+\frac{\partial f}{\partial u}\frac{\partial u}{\partial x_i}+\frac{\partial f}{\partial v}\frac{\partial v}{\partial x_i} \ \text{and}\  \frac{\partial f}{\partial y_i}=\frac{\partial f}{\partial s}\frac{\partial s}{\partial y_i}+\frac{\partial f}{\partial u}\frac{\partial u}{\partial y_i}+\frac{\partial f}{\partial v}\frac{\partial v}{\partial y_i}\  \text{for}\  i=1,2,3,4.
       \end{eqnarray*}
This gives the following equations:
\begin{equation}{\label{lemsimp vi}}
\frac{\partial f}{\partial x_1}=\frac{\partial f}{\partial s}\frac{\partial s}{\partial x_1}+\frac{\partial f}{\partial u}\frac{\partial u}{\partial x_1}=\frac{\partial \Tilde{\psi}_r}{\partial s}+\frac{\partial \Tilde{\psi}_r}{\partial u}y_1 \implies \frac{\partial f}{\partial x_1}|_{\alpha_0}=\frac{\partial \Tilde{\psi}_r}{\partial s}|_{\beta_0}+e\frac{\partial \Tilde{\psi}_r}{\partial u}|_{\beta_0}. 
\end{equation} 
\begin{equation}{\label{lemsimp vii}}
\frac{\partial f}{\partial x_2}=\frac{\partial \Tilde{\psi}_r}{\partial u}y_3 \implies \frac{\partial f}{\partial x_2}|_{\alpha_0}=m\frac{\partial \Tilde{\psi}_r}{\partial u}|_{\beta_0}.
       \end{equation}
\begin{equation}{\label{lemsimp viii}}
\frac{\partial f}{\partial x_3}=\frac{\partial \Tilde{\psi}_r}{\partial u}y_2\implies \frac{\partial f}{\partial x_3}|_{\alpha_0}=n\frac{\partial \Tilde{\psi}_r}{\partial u}|_{\beta_0}.
\end{equation} 
\begin{equation}{\label{lemsimp ix}}
      \frac{\partial f}{\partial x_4}=\frac{\partial \Tilde{\psi}_r}{\partial s}+\frac{\partial \Tilde{\psi}_r}{\partial u}y_4\implies \frac{\partial f}{\partial x_4}|_{\alpha_0}=\frac{\partial \Tilde{\psi}_r}{\partial s}|_{\beta_0}+w\frac{\partial \Tilde{\psi}_r}{\partial u}|_{\beta_0}.
      \end{equation}
       \begin{equation}{\label{lemsimp x}}
       \frac{\partial f}{\partial y_1}=\frac{\partial f}{\partial v}\frac{\partial v}{\partial y_1}+\frac{\partial f}{\partial u}\frac{\partial u}{\partial y_1}=\frac{\partial \Tilde{\psi}_r}{\partial v}+\frac{\partial \Tilde{\psi}_r}{\partial u}x_1\implies \frac{\partial f}{\partial y_1}|_{\alpha_0}=\frac{\partial \Tilde{\psi}_r}{\partial v}|_{\beta_0}+a\frac{\partial \Tilde{\psi}_r}{\partial u}|_{\beta_0}
       \end{equation}
       \begin{equation}{\label{lemsimp xi}}
       \frac{\partial f}{\partial y_2}=\frac{\partial \Tilde{\psi}_r}{\partial u}x_3\implies \frac{\partial f}{\partial y_2}|_{\alpha_0}=c\frac{\partial \Tilde{\psi}_r}{\partial u}|_{\beta_0}. 
       \end{equation}
       \begin{equation}{\label{lemsimp xii}}
       \frac{\partial f}{\partial y_3}=\frac{\partial \Tilde{\psi}_r}{\partial u}x_2\implies \frac{\partial f}{\partial y_3}|_{\alpha_0}=b\frac{\partial \Tilde{\psi}_r}{\partial u}|_{\beta_0}.
       \end{equation}
       \begin{equation}{\label{lemsimp xiii}}
       \frac{\partial f}{\partial y_4}=\frac{\partial \Tilde{\psi}_r}{\partial v}+\frac{\partial \Tilde{\psi}_r}{\partial u}x_4\implies\frac{\partial f}{\partial y_4}|_{\alpha_0}=\frac{\partial \Tilde{\psi}_r}{\partial v}|_{\beta_0}+d\frac{\partial \Tilde{\psi}_r}{\partial u}|_{\beta_0}. 
       \end{equation}
Now, $\left[\begin{array}{cccccccc}
\frac{\partial f}{\partial x_1} & \frac{\partial f}{\partial x_2} & \frac{\partial f}{\partial x_3} & \frac{\partial f}{\partial x_4} & \frac{\partial f}{\partial y_1} & \frac{\partial f}{\partial y_2} & \frac{\partial f}{\partial y_3} & \frac{\partial f}{\partial y_4}
\end{array}\right]\Big|_{\alpha_0} = \mathbf{0}$ implies \begin{enumerate}
\item  $(e-w)\frac{\partial\Tilde{\psi}_r}{\partial u}|_{\beta_0}=0$ by  Equation~(\ref{lemsimp vi}) and Equation~(\ref{lemsimp ix}).
\item $(a-d)\frac{\partial\Tilde{\psi}_r}{\partial u}|_{\beta_0}=0$ by  Equation~(\ref{lemsimp x}) and Equation~(\ref{lemsimp xiii}).
\item $m\frac{\partial \Tilde{\psi}_r}{\partial u}|_{\beta_0}=0=n\frac{\partial \Tilde{\psi}_r}{\partial u}|_{\beta_0}$ by  Equation~(\ref{lemsimp vii}) and Equation~(\ref{lemsimp viii}).
\item $c\frac{\partial \Tilde{\psi}_r}{\partial u}|_{\beta_0}=0=b\frac{\partial \Tilde{\psi}_r}{\partial u}|_{\beta_0}$ by  Equation~(\ref{lemsimp xi}) and Equation~(\ref{lemsimp xii}).
\end{enumerate}
If possible let $\frac{\partial \Tilde{\psi}_r}{\partial u}|_{\beta_0}\neq 0$. Then $A_0=aI$ and $B_0=eI$. As $A_0,B_0$ are in $\SL_2(k)$ therefore these conditions give $a=\pm 1$ and $e=\pm 1$. Therefore $\Bar{e}_2(A_0,B_0)=I$ and $\lambda_0=2$ which is a contradiction. Therefore  $\frac{\partial \Tilde{\psi}_r}{\partial u}|_{\beta_0}= 0$. 
       
This and the condition
$$\left[\begin{array}{cccccccc}
\frac{\partial f}{\partial x_1} & \frac{\partial f}{\partial x_2} & \frac{\partial f}{\partial x_3} & \frac{\partial f}{\partial x_4} & \frac{\partial f}{\partial y_1} & \frac{\partial f}{\partial y_2} & \frac{\partial f}{\partial y_3} & \frac{\partial f}{\partial y_4}
\end{array}\right]\Big|_{\alpha_0} = \mathbf{0}$$  together with the Equation~(\ref{lemsimp vi}) and Equation~(\ref{lemsimp x}), implies that $\left[\begin{array}{cccccccc}
\frac{\partial\Tilde{\psi}_r}{\partial s} & \frac{\partial\Tilde{\psi}_r}{\partial u} & \frac{\partial\Tilde{\psi}_r}{\partial v}\end{array}\right]\Big|_{\beta_0}=\mathbf{0}$. Hence the lemma.
  \end{proof}

\end{document}
